\documentclass[11pt,a4paper]{article}

% ---------------------------------------------------------------------------
% Packages
% ---------------------------------------------------------------------------
\usepackage[T1]{fontenc}
\usepackage[utf8]{inputenc}
\usepackage{lmodern}
\usepackage{microtype}
\usepackage[margin=2.4cm,headheight=14pt]{geometry}
\usepackage[table]{xcolor}
\usepackage{amsmath}
\usepackage{booktabs}
\usepackage{longtable}
\usepackage{array}
\usepackage{enumitem}
\usepackage{listings}
\usepackage{fancyhdr}
\usepackage{titlesec}
\usepackage[hidelinks,bookmarksnumbered]{hyperref}

% ---------------------------------------------------------------------------
% The cube palette (identical to DEFAULT_PALETTE in the engine)
% ---------------------------------------------------------------------------
\definecolor{white}{RGB}{255,255,255}
\definecolor{red}{RGB}{196,30,58}
\definecolor{green}{RGB}{0,158,96}
\definecolor{yellow}{RGB}{255,213,0}
\definecolor{orange}{RGB}{255,88,0}
\definecolor{blue}{RGB}{0,81,186}
\definecolor{codebg}{RGB}{247,247,247}
\definecolor{rulegray}{RGB}{170,170,170}

% One cell of a cube net: a colored square, or a blank of the same size.
\newcommand{\cubecell}[1]{\cellcolor{#1}\rule{4.6mm}{4.6mm}}
\newcommand{\cubeblank}{\phantom{\rule{4.6mm}{4.6mm}}}
\newcommand{\netfig}[1]{%
  \setlength{\arrayrulewidth}{0.3pt}%
  \renewcommand{\arraystretch}{0}%
  \begin{tabular}{@{}*{12}{c}@{}}\input{#1}\end{tabular}}

% ---------------------------------------------------------------------------
% Listings
% ---------------------------------------------------------------------------
\lstdefinestyle{shell}{
  basicstyle=\ttfamily\small,
  backgroundcolor=\color{codebg},
  frame=single,
  rulecolor=\color{rulegray},
  columns=fullflexible,
  keepspaces=true,
  breaklines=true,
  postbreak=\mbox{\textcolor{rulegray}{$\hookrightarrow$}\space},
  showstringspaces=false,
  aboveskip=0.8em,
  belowskip=0.8em
}
\lstdefinestyle{python}{
  style=shell,
  language=Python,
  keywordstyle=\bfseries,
  commentstyle=\itshape\color{black!55},
  morekeywords={as,with,match,case}
}
\lstdefinestyle{data}{
  style=shell,
  morestring=[b]"
}
\lstset{style=shell}

% ---------------------------------------------------------------------------
% Page style
% ---------------------------------------------------------------------------
\pagestyle{fancy}
\fancyhf{}
\fancyhead[L]{\small rubiks-cfop}
\fancyhead[R]{\small\thepage}
\renewcommand{\headrulewidth}{0.3pt}

\titleformat{\section}{\normalfont\Large\bfseries}{\thesection}{0.7em}{}
\titleformat{\subsection}{\normalfont\large\bfseries}{\thesubsection}{0.6em}{}

% ---------------------------------------------------------------------------
% Document
% ---------------------------------------------------------------------------
\title{\textbf{rubiks-cfop}\\[0.3em]
  \large A user manual for the WCA scramble generator and CFOP solver}
\author{The rubiks-cfop project}
\date{Version 0.1.0}

\begin{document}

\maketitle

\begin{abstract}
\noindent
\texttt{rubiks-cfop} is a headless engine, and a small command line program built
on it, that generates WCA scrambles for the $3\times3\times3$ Rubik's Cube and
solves them with the CFOP method.  The cube is represented with Kociemba's facelet
notation, every standard and wide move is implemented, and the known algorithms
live in JSON lookup tables.  This manual explains how to install and use the
program, how to read its output, how the solver works, and how to drive the engine
from Python.  Appendix~\ref{app:gui} sets out a plan for the graphical interface
that will be built on the same engine.
\end{abstract}

\tableofcontents
\clearpage

\section{Introduction}
\label{sec:intro}

\subsection{What the program does}

\texttt{rubiks-cfop} does two things:

\begin{enumerate}[leftmargin=1.4em]
  \item it \textbf{generates} scrambles for the $3\times3\times3$ cube in the form
        the World Cube Association uses, and
  \item it \textbf{solves} a cube with \textbf{CFOP}: \emph{Cross}, \emph{F2L}
        (the first two layers, four corner--edge pairs), \emph{OLL} (orientation of
        the last layer) and \emph{PLL} (permutation of the last layer).
\end{enumerate}

The program is two layers thick.  A \emph{headless engine} does all of the work and
returns structured data: it never prints, never reads the environment and has no
third-party runtime dependencies.  A thin \emph{command line interface} turns
arguments into engine calls and data into text.  A graphical interface can be
added later without touching the engine, which is the subject of
Appendix~\ref{app:gui}.

\subsection{Conventions}

\begin{description}[leftmargin=0em,style=nextline]
  \item[Cube representation] Kociemba's facelet notation: 54 facelets ordered
  \texttt{U1..U9}, \texttt{R1..R9}, \texttt{F1..F9}, \texttt{D1..D9},
  \texttt{L1..L9}, \texttt{B1..B9}.  The solved cube is the string
  \texttt{UUUUUUUUURRRRRRRRRFFFFFFFFFDDDDDDDDDLLLLLLLLLBBBBBBBBB}.
  \item[Moves] WCA notation (Regulation 12a): the face turns
  \texttt{U D L R F B}, the wide turns \texttt{Uw Dw Lw Rw Fw Bw}, the rotations
  \texttt{x y z}, and --- as a documented extension --- the slice turns
  \texttt{M E S}.  Every move accepts the suffixes \texttt{2} and \texttt{'}.
  \item[Colors] the standard scheme of Regulation 4d1, white on top and green in
  front: \texttt{U} white, \texttt{R} red, \texttt{F} green, \texttt{D} yellow,
  \texttt{L} orange, \texttt{B} blue.
  \item[Cross] always solved into the \textbf{bottom layer}.  For a white cross the
  upper layer is yellow, and the cube is reoriented so that this is what you see.
\end{description}

\subsection{WCA compliance}

The scramble generator follows Regulation 4b3, which asks an official scramble
program to ``produce a random state from all states that require at least 2 moves
to solve (equal probability for each state)''.  For the $3\times3\times3$ cube that
means \emph{random-state} scrambles: a uniformly random cube state is drawn and then
solved, and the inverse of that solution is the scramble.  Random-move sequences ---
fast, filtered, and what many tools mean by ``a scramble'' --- are also available,
but they are not WCA-legal for this puzzle (Regulation 4b3e permits them only for
$5\times5\times5$ and larger).  Scrambling orientation follows Regulation 4d1.

\subsection{Where the algorithms come from}

The OLL (57 cases) and PLL (21 cases) and F2L (41 cases) algorithms in the JSON
tables were taken from the published \emph{Logiqx cubing-algs} collection, whose
case identifiers match the community numbering.  They are not trusted on the
engineer's word: every entry is \emph{verified by execution}, and the recognizer's
patterns are \emph{derived from the algorithms themselves}, so the table and the
recognition can never disagree.  Section~\ref{sec:tables} describes the mechanism.

\subsection{How this manual is organized}

Section~\ref{sec:install} installs the program, Section~\ref{sec:quickstart} shows
it working, Section~\ref{sec:cli} is the complete command line reference and
Section~\ref{sec:output} explains the output line by line.
Section~\ref{sec:notation} is a notation primer, Section~\ref{sec:solver} explains
how each stage is solved, Section~\ref{sec:tables} documents the JSON tables and
Section~\ref{sec:library} the Python API.  Section~\ref{sec:trouble} covers
troubleshooting and the known limitations, and Appendix~\ref{app:gui} plans the
graphical interface.

\section{Installation}
\label{sec:install}

\subsection{Requirements}

\begin{itemize}[leftmargin=1.4em]
  \item \textbf{Python 3.13} or newer.  The project is pinned to 3.13 and the
        test suite asserts the interpreter version.
  \item \textbf{uv}, which manages the interpreter, the virtual environment and the
        development tools.
  \item No runtime dependency: the engine uses the standard library only.
\end{itemize}

\subsection{From a checkout}

\begin{lstlisting}[style=shell]
# fetch a managed Python 3.13 if you do not have one
uv python install 3.13

# create .venv, install the project and the development tools
uv sync

# the stage gate: formatting, linting and the whole test suite
scripts/check.sh
\end{lstlisting}

\texttt{uv sync} installs the project in editable mode and puts the
\texttt{rubiks-cfop} console script in \texttt{.venv/bin}.  You can run it directly,
through \texttt{uv run}, or through \texttt{uv tool install .} if you prefer a
standalone command.

\begin{lstlisting}[style=shell]
uv run rubiks-cfop --scramble "R U R' U' F2"
\end{lstlisting}

\subsection{Building a wheel}

\begin{lstlisting}[style=shell]
uv build            # dist/rubiks_cfop-0.1.0-py3-none-any.whl and .tar.gz
\end{lstlisting}

The wheel is self-contained: the JSON lookup tables live inside the package
(\texttt{rubiks\_cfop/data/}) and are installed with it, and the console script is
declared as an entry point.

\subsection{What is installed}

\begin{center}
\begin{tabular}{@{}ll@{}}
\toprule
Path & Contents \\
\midrule
\texttt{rubiks\_cfop/cube/} & facelets, moves, notation, \texttt{Cube}, \texttt{CubieCube} \\
\texttt{rubiks\_cfop/scramble/} & random-state and random-move generators \\
\texttt{rubiks\_cfop/solver/} & cross, F2L, OLL/PLL, CFOP pipeline, two-phase \\
\texttt{rubiks\_cfop/render/} & the unfolded net and its text serializers \\
\texttt{rubiks\_cfop/data/} & \texttt{oll.json}, \texttt{pll.json}, \texttt{f2l.json} \\
\texttt{rubiks\_cfop/cli.py} & the \texttt{rubiks-cfop} command \\
\bottomrule
\end{tabular}
\end{center}

\subsection{The first run is slower}

The random-state scrambler needs a set of precomputed tables: coordinates, move
tables and four pruning tables totalling about six megabytes.  They are built on
first use and then cached for the life of the process, which takes roughly ten
seconds of CPU.  Anything that runs the scrambler should build them up front on a
background thread:

\begin{lstlisting}[style=python]
import rubiks_cfop as cube

cube.preload()          # build the scrambler tables now
scramble = cube.generate_scramble(seed=1)
\end{lstlisting}

Solving a cube you already have does not need those tables, only the small cross
table (one megabyte, under a second).

\section{Quick start}
\label{sec:quickstart}

\subsection{Solve a scramble you were given}

\begin{lstlisting}[style=shell]
uv run rubiks-cfop --scramble "R U R' U' F2 D"
\end{lstlisting}

The program prints the scrambled cube as an unfolded net in real colors, then the
solution stage by stage.  Piped into a file, or with \texttt{--no-color}, the net is
drawn with colored squares instead of terminal colors.  A run looks like this:

\begin{lstlisting}[style=shell]
           Top
        ### ### ###
        ### ### ###    <- U face (yellow when the cross is white)
        ### ### ###
  Left    Front    Right    Back
### ###  ### ###  ### ###  ### ###  ...
### ###  ### ###  ### ###  ### ###  ...   <- L, F, R and B in a band
### ###  ### ###  ### ###  ### ###  ...
           Down
        ### ### ###
        ### ### ###    <- D face (the cross color)
        ### ### ###

Orientation: none
Cross color: yellow (bottom layer)
Scramble: R U R' U' F2 D

Cross: D' F2
F2L1: R' F R F'
F2L2: (none)
F2L3: (none)
F2L4: (none)
OLL (37): F R' F' R U R U' R'
PLL (skip): (none)

Total: 14 moves
\end{lstlisting}

(The \texttt{\#} characters stand for colored cells; Section~\ref{sec:output}
shows the real thing.)

\subsection{Generate a scramble and solve it}

\begin{lstlisting}[style=shell]
uv run rubiks-cfop --random --seed 42
\end{lstlisting}

\texttt{--random} is the default when no scramble is given.  The scramble is a
random-state scramble (WCA 4b3), typically 20 to 22 moves, and \texttt{--seed}
makes it reproducible: the same seed always yields the same scramble.

\begin{lstlisting}[style=shell]
Orientation: zz
Cross color: white (bottom layer)
Scramble: D B2 F2 L2 U' F2 U B2 R2 U L2 F R' U' B' U L R F D B

Cross: z z R' F D' F2 R' F
F2L1: U2 B' R2 B R2
F2L2: F' L F L2 U' L
F2L3: U L U' B U B' L'
F2L4: U B U B' U2 R' U2 R
OLL (40): U R' F R U R' U' F' U R
PLL (F Perm): U' L R2 U R U R2 U' R' U' R2 U' R U2 L' U R' U2

Total: 62 moves
\end{lstlisting}

\subsection{Pick the cross yourself}

\begin{lstlisting}[style=shell]
uv run rubiks-cfop --scramble "R U R' U'" --cross white
uv run rubiks-cfop --scramble "R U R' U'" --cross-options
\end{lstlisting}

\texttt{--cross} names one of the six colors; \texttt{--cross-options} first lists
what each of the six crosses would cost, shortest first:

\begin{lstlisting}[style=shell]
  yellow  0 moves
  white   3 moves
  green   4 moves
  ...
\end{lstlisting}

\subsection{Get the result as JSON}

\begin{lstlisting}[style=shell]
uv run rubiks-cfop --scramble "R U R' U'" --json
\end{lstlisting}

The payload carries the scramble, the cross color, the rotation applied, every
stage with its moves and the facelets \emph{after} the stage, the metrics and a
verification flag.  It is the same structure the Python API returns from
\texttt{Solution.to\_dict()}, so a GUI or a web front-end can consume it directly;
Section~\ref{sec:library} documents it.

\section{Command line reference}
\label{sec:cli}

\begin{lstlisting}[style=shell]
rubiks-cfop [--scramble MOVES | --facelets STATE | --random]
            [--scramble-style {random-state,random-move}] [--seed N]
            [--cross {best,white,red,green,yellow,orange,blue}]
            [--color | --no-color] [--json] [--cross-options]
\end{lstlisting}

\subsection{Where the cube comes from}

Exactly one of these three sources is used.  With none of them the program behaves
as if \texttt{--random} had been given.

\begin{longtable}{@{}p{3.2cm}p{11.2cm}@{}}
\toprule
Option & Meaning \\
\midrule
\endhead
\texttt{--scramble MOVES} &
Apply a scramble in WCA notation to a solved cube.  Wide moves
(\texttt{Rw}, \texttt{2Rw}), rotations (\texttt{x y z}) and slice moves
(\texttt{M E S}) are accepted, as are the usual aliases (\texttt{r} for
\texttt{Rw}, \texttt{m} for \texttt{M}).  Invalid notation is reported with the
offending token and its position, and the program exits with status 2. \\
\texttt{--facelets STATE} &
Start from a 54 character Kociemba facelet string instead of a scramble.  Useful
when another tool produced the state, or when you want to solve a state you built
by hand.  Whitespace and lower case are ignored; impossible states (a twisted
corner, a wrong colour count, moved centres) are rejected. \\
\texttt{--random} &
Generate a scramble and solve it.  This is the default. \\
\bottomrule
\end{longtable}

\subsection{How the scramble is generated}

\begin{longtable}{@{}p{3.2cm}p{11.2cm}@{}}
\toprule
Option & Meaning \\
\midrule
\endhead
\texttt{--scramble-style} &
\texttt{random-state} (default) draws a uniformly random cube state and solves it,
which is what WCA Regulation 4b3 requires; scrambles are 20 to 22 moves.
\texttt{random-move} produces a filtered random-move sequence of 20 face turns
instead --- instant, but \emph{not} WCA-legal for the $3\times3\times3$. \\
\texttt{--seed N} &
Seed for a generated scramble.  The same seed always gives the same scramble, and
the seed is echoed in the JSON payload.  Without a seed one is drawn from the
operating system and reported, so any run can be reproduced. \\
\bottomrule
\end{longtable}

\subsection{How it is solved}

\begin{longtable}{@{}p{3.2cm}p{11.2cm}@{}}
\toprule
Option & Meaning \\
\midrule
\endhead
\texttt{--cross} &
The colour of the cross: \texttt{white}, \texttt{yellow}, \texttt{green},
\texttt{blue}, \texttt{red} or \texttt{orange}.  \texttt{best} (the default) solves
all six crosses and keeps the shortest, breaking ties in the order white, yellow,
green, blue, red, orange.  The cube is always reoriented so that the chosen colour
ends up in the \emph{bottom} layer. \\
\texttt{--cross-options} &
Before solving, print the six crosses with their lengths, shortest first. \\
\bottomrule
\end{longtable}

\subsection{How it is presented}

\begin{longtable}{@{}p{3.2cm}p{11.2cm}@{}}
\toprule
Option & Meaning \\
\midrule
\endhead
\texttt{--color} &
Force ANSI truecolour output, using 24-bit background colours and the standard
cubing palette. \\
\texttt{--no-color} &
Draw the net with coloured square characters instead, so the output survives being
redirected to a file or pasted into a document.  The default is automatic: colours
when standard output is a terminal, squares otherwise. \\
\texttt{--json} &
Print the machine-readable payload instead of the human listing. \\
\texttt{-h}, \texttt{--help} &
Usage summary. \\
\bottomrule
\end{longtable}

\subsection{Exit codes}

\begin{center}
\begin{tabular}{@{}cl@{}}
\toprule
Code & Meaning \\
\midrule
0 & A solution was produced (and verified). \\
2 & The input could not be understood --- bad notation, an impossible state --- or
the cube could not be solved.  The reason is written to standard error, prefixed
with \texttt{error:}. \\
\bottomrule
\end{tabular}
\end{center}

\subsection{Examples}

\begin{lstlisting}[style=shell]
# a scramble from a competition
uv run rubiks-cfop --scramble "D B2 F2 L2 U' F2 U B2 R2 U"

# reproducible random-state scramble, plain output for a log
uv run rubiks-cfop --random --seed 7 --no-color

# solve a state produced elsewhere, with a white cross of my choosing
uv run rubiks-cfop --facelets UULUUFUUFRRUBRRURRFFDFFUFFFDDRDDDDDDBLLLLLLLLBRRBBBBBB \
                   --cross white

# machine readable, for a script
uv run rubiks-cfop --scramble "R U R' U'" --json | jq .total_moves
\end{lstlisting}

\section{Reading the output}
\label{sec:output}

\subsection{The unfolded cube}

The first thing printed is the scrambled cube, unfolded into the usual cross shape:
the front face in the middle, the upper face above it, the bottom face below it and
the left, front, right and back faces across a band.

\begin{figure}[htbp]
\centering
\netfig{figures/net-solved}
\caption{The solved cube.  Top to bottom: U, then the band L--F--R--B, then D.}
\label{fig:net-solved}
\end{figure}

\begin{figure}[htbp]
\centering
\netfig{figures/net-scrambled}
\caption{The cube after the scramble \texttt{R U R' U' F2}.  Each square is one
facelet; Section~\ref{sec:notation} gives the facelet numbering.}
\label{fig:net-scrambled}
\end{figure}

The cells are drawn in the cube's real colours --- no letters describe a colour, so
the picture can be checked at a glance against the physical cube.

\begin{center}
\begin{tabular}{@{}llll@{}}
\toprule
Face & Colour & RGB & Position in the net \\
\midrule
\texttt{U} & white  & \texttt{\#FFFFFF} & above the band \\
\texttt{L} & orange & \texttt{\#FF5800} & band, first group \\
\texttt{F} & green  & \texttt{\#009E60} & band, second group \\
\texttt{R} & red    & \texttt{\#C41E3A} & band, third group \\
\texttt{B} & blue   & \texttt{\#0051BA} & band, fourth group \\
\texttt{D} & yellow & \texttt{\#FFD500} & below the band \\
\bottomrule
\end{tabular}
\end{center}

\subsection{The orientation lines}

\begin{lstlisting}[style=shell]
Orientation: zz
Cross colour: white (bottom layer)
\end{lstlisting}

The cross is always solved into the bottom layer, so the cube you are shown is
usually \emph{not} in the orientation it was scrambled in.  \texttt{Orientation}
records the whole-cube rotation that was applied to put the chosen colour at the
bottom (\texttt{none} when no rotation was needed).  For a white cross the top face
is therefore yellow, as the second line states.

Both lines describe the same picture: compare the bottom face of the net with the
\texttt{Cross colour} line, and the top face with its opposite.

\subsection{The stages}

\begin{lstlisting}[style=shell]
Cross: D' F2
F2L1: R' F R F'
F2L2: (none)
F2L3: (none)
F2L4: (none)
OLL (37): F R' F' R U R U' R'
PLL (skip): (none)

Total: 14 moves
\end{lstlisting}

\begin{description}[leftmargin=0em,style=nextline]
  \item[\texttt{Cross}] the moves that solve the four bottom edges.  Optimal: no
  cross needs more than eight moves, and this one is as short as possible for the
  chosen colour.
  \item[\texttt{F2L1} to \texttt{F2L4}] one line per inserted pair, in the order
  front-right, front-left, back-left, back-right.  Each line keeps the cross and
  every earlier pair solved.
  \item[\texttt{OLL (\textit{n})}] the orientation of the last layer, with the
  number of the case in the standard 1--57 numbering.  When the last layer is
  already oriented the line reads \texttt{OLL (skip)} and has no moves.
  \item[\texttt{PLL (\textit{name})}] the permutation of the last layer, with the
  name of the case (\texttt{T Perm}, \texttt{Ua Perm}, \ldots).  When the last layer
  is already permuted the line reads \texttt{PLL (skip)}.
  \item[\texttt{(none)}] a stage that needed no moves.
  \item[\texttt{Total}] the number of moves in the whole solution.
\end{description}

A stage's move list occasionally contains a rotation (\texttt{x}, \texttt{y} or
\texttt{z}) and wide or slice moves: some published algorithms use them, and the
rotations keep the cube in the ``cross down'' orientation the net is drawn in.

\subsection{Applying the solution}

Read the stages top to bottom and apply the moves in the order printed to the cube
in the orientation of the net --- the cross colour at the bottom.  The line
\texttt{Orientation} tells you how the cube was turned to get there, so you can
reproduce that orientation on a physical cube.  Every move is standard WCA notation
(Section~\ref{sec:notation}).

A quick check that you are reading it right: after \texttt{Cross}, the four bottom
edges and their side centres must form a plus sign, and the bottom face must be
entirely the cross colour.

\subsection{The JSON payload}

\begin{lstlisting}[style=data]
{
  "scramble": "R U R' U'",
  "cross_color": "yellow",
  "rotation": "",
  "stages": [
    {
      "label": "Cross", "kind": "cross", "moves": [], "text": "", "length": 0,
      "facelets": "UULUUFUUFRRUBRRURRFFDFFUFFFDDRDDDDDDBLLLLLLLLBRRBBBBBB",
      "metadata": {"color": "yellow", "rotation": ""}
    },
    {
      "label": "F2L1", "kind": "f2l", "moves": ["R'", "F", "R", "F'"],
      "text": "R' F R F'", "length": 4,
      "facelets": "UUFUUFBBURUURRRRRRUUFFFFFFFDDDDDDDDDBLLLLLLLLLRRBBBBBB",
      "metadata": {"slot": "FR", "case": "1", "case_name": "Same Colour on Top"}
    },
    ...
  ],
  "total_moves": 12,
  "verified": true,
  "final_facelets": "UUUUUUUUURRRRRRRRRFFFFFFFFFDDDDDDDDDLLLLLLLLLBBBBBBBBB"
}
\end{lstlisting}

\texttt{facelets} on each stage is the state \emph{after} that stage, so a user
interface can show or animate any stage without re-running the solver, and
\texttt{verified} is the result of re-applying the whole solution to the scramble.

\section{Notation}
\label{sec:notation}

\subsection{Face turns}

A letter names a face and the letter alone means a quarter turn clockwise as seen
from outside that face.  The suffix \texttt{2} means a half turn and \texttt{'}
means a quarter turn anticlockwise (WCA Regulation 12a1).

\begin{center}
\begin{tabular}{@{}llll@{}}
\toprule
Letter & Face & Clockwise seen from & Opposite \\
\midrule
\texttt{U} & upper        & above   & \texttt{D} \\
\texttt{D} & bottom       & below   & \texttt{U} \\
\texttt{R} & right        & right   & \texttt{L} \\
\texttt{L} & left         & left    & \texttt{R} \\
\texttt{F} & front        & front   & \texttt{B} \\
\texttt{B} & back         & behind  & \texttt{F} \\
\bottomrule
\end{tabular}
\end{center}

There are 18 face turns: \texttt{U U2 U' R R2 R' F F2 F' D D2 D' L L2 L' B B2 B'}.

\subsection{Wide turns}

An \texttt{w} after the letter turns the outer face \emph{and} the layer behind it
(Regulation 12a2): \texttt{Uw Dw Lw Rw Fw Bw} with the usual suffixes, 18 moves in
total.  The layer count may be written out, so \texttt{2Rw} is the same move as
\texttt{Rw}; \texttt{1Rw} is not valid notation and \texttt{3Rw} does not exist on a
$3\times3\times3$.

\subsection{Rotations and slices}

A rotation turns the \emph{whole} cube, including its centres, in the direction of
a face turn (Regulation 12a4): \texttt{x} follows \texttt{R}, \texttt{y} follows
\texttt{U} and \texttt{z} follows \texttt{F}, with the usual suffixes.

The slice turns \texttt{M}, \texttt{E} and \texttt{S} turn only the middle layer,
following \texttt{L}, \texttt{D} and \texttt{F} respectively.  They are not part of
WCA notation but appear in many published algorithms, so the engine understands
them; Section~\ref{sec:tables} explains how the CLI keeps the display upright when
an algorithm uses them.

\subsection{Accepted aliases}

\begin{center}
\begin{tabular}{@{}ll@{}}
\toprule
Written & Read as \\
\midrule
\texttt{r l u d f b} & \texttt{Rw Lw Uw Dw Fw Bw} (older listings) \\
\texttt{m e s} & \texttt{M E S} \\
\texttt{RW} & \texttt{Rw} (the \texttt{w} is case insensitive) \\
\texttt{2Rw} & \texttt{Rw} \\
\texttt{R2'} or \texttt{R'2} & \texttt{R2} \\
\bottomrule
\end{tabular}
\end{center}

Anything else is rejected with the offending token and its position, for example
\texttt{error: layer prefix '3' is invalid on a 3x3x3 cube (index=4, token='3Rw')}.

\subsection{The facelet representation}
\label{sec:facelets}

The cube is a string of 54 colours in Kociemba's order: the nine facelets of
\texttt{U}, then of \texttt{R}, \texttt{F}, \texttt{D}, \texttt{L} and \texttt{B}.
Inside a face the facelets are numbered row by row in reading order \emph{as seen
from outside the cube}, with the following view directions, which are the ones
Kociemba's own move definitions use:

\begin{center}
\begin{tabular}{@{}lllll@{}}
\toprule
Face & View up & Top-left facelet belongs to & Top-right & Bottom-right \\
\midrule
\texttt{U} & back          & UBL & UBR & UFR \\
\texttt{R} & up            & UFR & UBR & DBR \\
\texttt{F} & up            & UFL & UFR & DFR \\
\texttt{D} & front         & DFL & DFR & DBR \\
\texttt{L} & up            & UBL & UFL & DFL \\
\texttt{B} & up            & UBR & UBL & DBL \\
\bottomrule
\end{tabular}
\end{center}

For example, \texttt{U9} is the upper sticker of the front-right corner,
\texttt{R1} is the right sticker of that same corner, and \texttt{F3} is its front
sticker, so the three facelets \texttt{(U9, R1, F3)} describe one and the same
corner piece.  The solved cube is

\begin{center}
\texttt{UUUUUUUUURRRRRRRRRFFFFFFFFFDDDDDDDDDLLLLLLLLLBBBBBBBBB}
\end{center}

and the six centre facelets --- \texttt{U5}, \texttt{R5}, \texttt{F5}, \texttt{D5},
\texttt{L5}, \texttt{B5} --- say which colour belongs to which face.  A state in
which the centres have moved is still a legal cube (it has simply been turned in
the hand); the program handles that by relabelling the colours before it reasons
about pieces.

\subsection{Moves as permutations}

Every move is stored as a permutation of the 54 facelets: for each facelet, where
its sticker goes.  The eighteen face turns were generated from the cube's geometry
and then checked, one by one, against the move tables published on
\texttt{kociemba.org}; wide turns are two-layer rotations, rotations are three-layer
rotations and slices are the single middle layer.  That is why a wide move or a
rotation can move a centre sticker and a face turn cannot.

\section{How it solves}
\label{sec:solver}

CFOP is four stages.  Each one is a separate module, each returns structured data,
and the pipeline verifies the whole solution before handing it back.

\subsection{The cross}

The four bottom edges are solved \emph{optimally}.  Only three things matter about
a cross: which four of the twelve edge positions the bottom edges occupy, which of
the four pieces is in which position, and how each is oriented.  That gives
\[
  \binom{12}{4} \cdot 4! \cdot 2^4 = 495 \cdot 24 \cdot 16 = 190{,}080
\]
states.  A breadth-first search over that space yields a distance table of one byte
per state, about a megabyte, and a solution is read off by walking downhill one move
at a time.  The table takes under a second to build, once per process.

No cross needs more than eight moves, and because the table stores exact distances,
the reported cross is as short as it possibly can be.  \texttt{--cross best} simply
repeats the search for all six colors and keeps the shortest.

\subsection{The first two layers}

The four corner--edge pairs are inserted one at a time, in the order front-right,
front-left, back-left, back-right, each insertion keeping the cross and every pair
already inserted.

Two mechanisms cooperate:

\begin{description}[leftmargin=0em,style=nextline]
  \item[The 41 case table] \texttt{data/f2l.json} holds the standard cases, each
  expressed in the front-right frame.  A slot is recognized by the position and
  orientation of its two pieces, and an algorithm is applied --- conjugated by a
  whole-cube turn when the slot is not the front-right one, so the display stays
  cross-down.
  \item[The verified search] Anything the table does not cover --- which includes
  any configuration where one of the pieces is stuck in another slot --- is inserted
  by an iterative deepening search over the cube, guided by three admissible
  heuristics: the distance to insert the pair, the distance to restore the cross and
  the distance to restore each pair already inserted.  The search only stops at a
  state where all of them are solved, so it cannot break an earlier stage.
\end{description}

A table result is always checked against the same goal before it is accepted; if it
would disturb something, the search takes over.  The pair subproblem itself is
small --- eight corner positions times three orientations times twelve edge
positions times two orientations, 576 states --- which is what makes the heuristic
cheap to evaluate at every node.

\subsection{Orientation and permutation of the last layer}

OLL and PLL are lookup tables of 57 and 21 cases.  Recognition does not use
hand-written pattern descriptions, which are easy to get subtly wrong.  Instead the
pattern of a case is \emph{computed from the case}: for each algorithm the engine
applies its inverse to a solved cube, reads off the pattern of the resulting state,
and stores that as the recognition key.  Table and recognizer therefore cannot
drift apart.

A case is only defined up to an adjustment of the top face, so every case is stored
for all sixteen combinations of a turn \emph{before} the algorithm and a turn
\emph{after} it --- the two-sided AUF of the cubing literature.  That is what makes
the tables complete.  The test suite proves it by enumerating the whole last layer:

\begin{center}
\begin{tabular}{@{}lll@{}}
\toprule
Space & Size & Verified \\
\midrule
Orientation states (OLL) & $3^3 \cdot 2^3 = 216$ & every state recognized and oriented \\
Permutation states (PLL) & $4! \cdot 4! / 2 = 288$ & every state recognized and permuted \\
\bottomrule
\end{tabular}
\end{center}

The 216 and 288 include the already-solved configuration, which needs no algorithm:
the pipeline reports \texttt{OLL (skip)} or \texttt{PLL (skip)} instead.

\subsection{Scrambles: a two-phase solver}

Random-state scrambles need a solver that can handle any state quickly, so the
project includes a Kociemba-style \emph{two-phase} solver.  Six coordinates describe
the cube compactly:

\begin{center}
\begin{tabular}{@{}llr@{}}
\toprule
Coordinate & Meaning & Values \\
\midrule
twist        & orientation of the eight corners & $3^7 = 2187$ \\
flip         & orientation of the twelve edges & $2^{11} = 2048$ \\
slice        & which positions hold the four slice edges & $\binom{12}{4} = 495$ \\
corner perm  & permutation of the corners & $8! = 40{,}320$ \\
edge perm    & permutation of the eight up/down edges & $8! = 40{,}320$ \\
slice perm   & permutation of the four slice edges & $4! = 24$ \\
\bottomrule
\end{tabular}
\end{center}

Phase~1 brings twist, flip and slice home using all eighteen face turns; phase~2
solves the rest using only half turns of R, L, F, B and any turn of U and D, which
cannot disturb what phase~1 fixed.  Both phases are iterative deepening searches
guided by breadth-first pruning tables (the worst phase-1 distances are 9, 9, 14
and 12).  A solve takes about 20 to 22 moves and, once the tables are built, a
fraction of a second.

To produce a scramble the generator draws a uniformly random solvable state, solves
it, and returns the \emph{inverse} of that solution --- so applying the scramble to a
solved cube really does produce the drawn state.  States that would need fewer than
two moves are redrawn, as Regulation 4b3 requires.

\subsection{Every solution is verified}

The pipeline concatenates the stages and applies the result to the scramble.  If the
cube is not solved, it raises an error rather than returning the solution: a
solution that does not check out is a bug, not a result.  The same flag is exposed as
\texttt{verified} in the JSON payload.

\subsection{Measured performance}

Figures from the development machine, for a random-state scramble of 20 to 22 moves:

\begin{center}
\begin{tabular}{@{}lr@{}}
\toprule
Step & Typical cost \\
\midrule
Build the scrambler tables (first use only) & 10 s \\
Generate a random-state scramble & 0.1--0.5 s \\
Cross, all six colors & under 1 ms \\
Four F2L pairs & 0.5--4 s (search-bound; see below) \\
OLL and PLL & under 5 ms \\
Whole solve & 0.7--4 s, 54 to 62 moves \\
\bottomrule
\end{tabular}
\end{center}

The F2L search dominates.  It is only used for configurations the 41-case table does
not cover, and it is correct for every one of them; a complete pair table would make
it unnecessary.  This is the single largest performance improvement available and is
listed again in Section~\ref{sec:trouble}.

\section{The lookup tables}
\label{sec:tables}

The known algorithms live in JSON, one file per stage, inside the package:

\begin{center}
\begin{tabular}{@{}llr@{}l@{}}
\toprule
File & Stage & Cases & Labels \\
\midrule
\texttt{data/oll.json} & orientation of the last layer & 57 & \texttt{"1"} \ldots{} \texttt{"57"} \\
\texttt{data/pll.json} & permutation of the last layer & 21 & \texttt{"Aa Perm"} \ldots{} \texttt{"Z Perm"} \\
\texttt{data/f2l.json} & the four corner--edge pairs & 41 & \texttt{"1"} \ldots{} \texttt{"41"} \\
\bottomrule
\end{tabular}
\end{center}

\subsection{Schema}

All three files share one shape:

\begin{lstlisting}[style=data]
{
  "schema": 1,
  "kind": "oll",
  "cases": [
    {
      "label": "1",
      "name": "No Edges - Runway",
      "algorithm": "R U2' R2' F R F' U2 R' F R F'"
    }
  ]
}
\end{lstlisting}

\begin{description}[leftmargin=0em,style=nextline]
  \item[\texttt{schema}] the format version, so a future change can be detected
  rather than misread;
  \item[\texttt{kind}] \texttt{oll}, \texttt{pll} or \texttt{f2l}, which tells the
  loader which recognizer to build;
  \item[\texttt{label}] what the user sees in the solution listing --- the OLL
  number, the PLL name, or the F2L case number;
  \item[\texttt{name}] a longer description, shown by interfaces and in the JSON
  metadata;
  \item[\texttt{algorithm}] one algorithm in WCA notation.
\end{description}

\subsection{Where the algorithms come from}

The tables are generated, not typed by hand, from the published \emph{Logiqx
cubing-algs} collection, whose case identifiers match the community numbering:

\begin{lstlisting}[style=shell]
uv run python tools/build_last_layer_tables.py oll.js pll.js f2l.js
\end{lstlisting}

For every case the tool considers all the algorithms listed for it, keeps the ones
that parse, checks that each one really solves its own case, and prefers the
shortest algorithm that uses face turns only.  Those preferences matter: a wide or
slice move shifts the centers, so an algorithm that uses one would leave the cube
in a different orientation than the net is drawn in.

\subsection{Why recognition cannot drift}

The loader derives each case's recognition key from its algorithm, by applying the
inverse of the algorithm to a solved cube and reading the pattern of the result.  A
consequence worth knowing: if you edit an \texttt{algorithm} and restart, both the
solution and the recognition change together, and there is no second place to update.

Two cases whose patterns collide are a hard error --- the loader raises
\texttt{InvalidStateError} naming both labels --- because a state with two possible
algorithms is ambiguous.

\subsection{Adding or changing a case}

\begin{enumerate}[leftmargin=1.4em]
  \item edit the JSON file: add an entry with a fresh label, a name and an algorithm;
  \item run the tests.  The per-case test applies every algorithm to its own case
  state and requires the stage to reach its goal, so a wrong algorithm fails loudly;
  \item run \texttt{scripts/check.sh} and commit.
\end{enumerate}

Tables are loaded once per process and cached, so a running program keeps the tables
it started with --- restart to pick up an edit.

\subsection{What the tables do not cover}

Two gaps are worth stating plainly:

\begin{itemize}[leftmargin=1.4em]
  \item of the 41 published F2L cases, 35 survive the ``face turns only''
  preference.  Nothing is lost: the other six, and every configuration the 41 cases
  do not describe at all, are handled by the verified search.
  \item the last layer can contain a state for which the stored algorithm does not
  finish the cube.  The pipeline notices, finishes the remainder with the two-phase
  solver and reports it honestly in the stage metadata.
\end{itemize}

\section{Using the engine from Python}
\label{sec:library}

The engine is a library first and a program second.  Everything the command line
does is available as a function, and everything it prints is available as data.

\subsection{The public surface}

\begin{longtable}{@{}p{4.6cm}p{9.8cm}@{}}
\toprule
Name & Purpose \\
\midrule
\endhead
\texttt{Cube} & an immutable cube state: \texttt{apply}, \texttt{apply\_moves},
\texttt{apply\_sequence}, \texttt{apply\_scramble}, \texttt{is\_solved},
\texttt{is\_solvable}, \texttt{cubie}, \texttt{to\_dict} \\
\texttt{CubieCube} & the same cube at the piece level: corner and edge permutations
with orientations, plus \texttt{validate} \\
\texttt{MOVES}, \texttt{get\_move}, \texttt{invert} & the 54 canonical moves and
their inverses \\
\texttt{parse\_moves}, \texttt{parse\_scramble}, \texttt{format\_moves},
\texttt{normalize\_sequence} & notation in and out \\
\texttt{generate\_scramble}, \texttt{generate\_scrambles} & scrambles, with a seed
and a style \\
\texttt{Scramble} & the movement list, its text, the state it produces, its seed \\
\texttt{solve\_cube} & the CFOP pipeline: the whole solution as a \texttt{Solution} \\
\texttt{cross\_options} & the six cross colours with their lengths, ranked \\
\texttt{Solution}, \texttt{Stage}, \texttt{CrossOption} & the result objects, each
with \texttt{to\_dict} \\
\texttt{solve}, \texttt{SearchLimits} & the two-phase solver on its own, with depth
and time limits \\
\texttt{CancelToken} & cooperative cancellation for a solve \\
\texttt{preload} & build the scrambler tables now \\
\texttt{net\_layout}, \texttt{render\_ansi}, \texttt{render\_emoji},
\texttt{render\_letters}, \texttt{NetLayout}, \texttt{DEFAULT\_PALETTE} & the
unfolded net, as data or as text \\
\texttt{RubiksError} and subclasses & structured errors, each with
\texttt{to\_dict} \\
\bottomrule
\end{longtable}

\subsection{Scramble and solve}

\begin{lstlisting}[style=python]
import rubiks_cfop as cube

scramble = cube.generate_scramble(seed=42)
print(scramble.text)                 # "D B2 F2 ..."
print(scramble.style.value)          # "random-state"
print(scramble.state.is_solvable())  # True

solution = cube.solve_cube(scramble.state, scramble=scramble.text)
for stage in solution.stages:
    print(f"{stage.label}: {stage.text}")
print(solution.total_moves, solution.verified)
\end{lstlisting}

\begin{lstlisting}[style=shell]
Cross: z z R' F D' F2 R' F
F2L1: U2 B' R2 B R2
F2L2: F' L F L2 U' L
F2L3: U L U' B U B' L'
F2L4: U B U B' U2 R' U2 R
OLL (40): U R' F R U R' U' F' U R
PLL (F Perm): U' L R2 U R U R2 U' R' U' R2 U' R U2 L' U R' U2
62 True
\end{lstlisting}

\subsection{Choosing the cross}

\begin{lstlisting}[style=python]
solution = cube.solve_cube(cube.Cube.solved().apply_scramble("R U R' U'"), cross="white")

for option in cube.cross_options(scramble.state):
    print(option.color.value, option.length)
\end{lstlisting}

The colour may be a \texttt{Color} or its name; \texttt{"best"} is the default.

\subsection{Drawing the cube}

\begin{lstlisting}[style=python]
layout = cube.net_layout(scramble.state)   # 9 rows x 12 columns of Color
print(cube.render_emoji(layout))           # for a terminal without colour
print(cube.render_ansi(layout))            # 24-bit ANSI background colours

for face, rows in layout.to_dict()["faces"].items():
    print(face, rows[0])                   # the top row of each face, by name
\end{lstlisting}

\texttt{NetLayout} holds colours, not escape codes: an interface draws the cube in
its own way from the same structure.  \texttt{render\_letters} exists for logs and
diffs, and is deliberately not the default, because the program's output is meant to
show real colours.

\subsection{Cancelling a solve}

A solve can take a few seconds, and an interactive program must be able to abandon
one.

\begin{lstlisting}[style=python]
import threading
import rubiks_cfop as cube

cancel = cube.CancelToken()

def work():
    try:
        solution = cube.solve_cube(scramble.state)
    except cube.CancelledError:
        print("abandoned")

thread = threading.Thread(target=work)
thread.start()
# ... later, from the user interface thread:
cancel.cancel()
\end{lstlisting}

Passing a token to a solve is a matter of scheduling: the pipeline checks it at
regular intervals and raises \texttt{CancelledError} promptly.  The two-phase solver
accepts the same token and an optional timeout through \texttt{SearchLimits}.

\subsection{Errors}

Every error derives from \texttt{RubiksError} and carries a stable \texttt{code}, a
message and structured details, so a caller never has to parse a string:

\begin{center}
\begin{tabular}{@{}ll@{} }
\toprule
Class & When \\
\midrule
\texttt{NotationError} & a move sequence that cannot be parsed, with
\texttt{token} and \texttt{index} \\
\texttt{InvalidMoveError} & a single move that is not valid on a $3\times3\times3$ \\
\texttt{InvalidStateError} & a malformed state: wrong length, wrong counts, moved
centres, a piece that cannot exist \\
\texttt{UnsolvableStateError} & a well formed but impossible state: a twisted
corner, a flipped edge, wrong parity \\
\texttt{SearchLimitError} & a search ran out of depth budget \\
\texttt{SolveTimeoutError} & a solve ran out of time \\
\texttt{CancelledError} & a solve was cancelled \\
\bottomrule
\end{tabular}
\end{center}

\begin{lstlisting}[style=python]
try:
    cube.Cube.solved().apply_sequence("R Q")
except cube.NotationError as error:
    print(error.to_dict())
    # {'code': 'notation-error', 'message': 'not WCA move notation',
    #  'details': {'index': 1, 'token': 'Q'}}
\end{lstlisting}

\subsection{Threading and immutability}

\texttt{Cube}, \texttt{CubieCube}, \texttt{Move}, \texttt{Solution} and
\texttt{Stage} are frozen value objects: sharing one between threads is safe, and
applying a move returns a new cube rather than changing one.  Solved data is never
mutated, the caches are read-only tables, and the solvers keep their state in local
variables, so a program may solve on a worker thread while its interface keeps
running.  The only expensive one-off work --- the scrambler tables --- is exposed as
\texttt{preload()} precisely so that it can be moved off the interactive path.

\section{Troubleshooting and known limitations}
\label{sec:trouble}

\subsection{Symptoms and what they mean}

\begin{longtable}{@{}p{5.4cm}p{9cm}@{}}
\toprule
Symptom & Cause and remedy \\
\midrule
\endhead
The first solve takes ten seconds & The scrambler tables are being built.  They are
cached for the life of the process; call \texttt{preload()} at start-up to pay the
cost once, on a thread of your own. \\
\texttt{error: not WCA move notation} & A token in the scramble is not a move.  The
message names the token and its index; \texttt{3Rw} and \texttt{1Rw} are invalid on
a $3\times3\times3$. \\
\texttt{error: the corner orientation sum must be divisible by three} & The state is
impossible: a single corner is twisted.  Re-check the facelets, or the tool that
produced them. \\
\texttt{error: each color must appear exactly nine times} & A facelet string with
nine of each color is the minimum requirement; a missing or duplicated sticker
usually means a transcription slip. \\
\texttt{error: the last layer does not match any case} & The last layer is not in a
state the tables know.  The pipeline normally catches this and finishes with the
two-phase solver; seeing the error directly means a table file is damaged. \\
\texttt{SearchLimitError} & A search was given too small a depth budget.  Raise
\texttt{SearchLimits.max\_total\_depth} (the default of 30 is generous) or let the
pair search run to its default depth of twelve. \\
\texttt{SolveTimeoutError} & A solve exceeded \texttt{SearchLimits.timeout\_ms} or
\texttt{SearchLimits.max\_nodes}.  Scramble generation gives each attempt a budget of
400,000 search nodes (about a second) and redraws a state that overruns; counting
nodes rather than time keeps a seed's scramble identical on every machine. \\
\texttt{CancelledError} & Someone called \texttt{CancelToken.cancel()}.  Expected
when an interface abandons a solve; reuse the token with \texttt{reset()}. \\
Solving feels slow & The F2L search is doing the work: the 41-case table only
covers configurations whose pieces are in the U layer or the slot.  Correct, but a
second or two per pair. \\
The net looks upside down & It is: the cross is at the bottom.  The
\texttt{Orientation} line says how the cube was turned to get there. \\
\bottomrule
\end{longtable}

\subsection{Known limitations}

\begin{enumerate}[leftmargin=1.4em]
  \item \textbf{F2L speed.}  The largest available improvement.  A complete table
  over the 576 pair configurations (or an explicit extraction step before the case
  table) would remove the search from the common path.
  \item \textbf{Last-layer naming.}  For a small number of states the stored
  algorithm does not finish the cube; the pipeline completes it with the two-phase
  solver instead of a named case.  The solution is still correct and verified, but
  the stage is reported as such rather than as \texttt{OLL (\textit{n})}.
  \item \textbf{Rotations inside algorithms.}  A few published algorithms include a
  rotation or a wide or slice move.  The stage re-orients the display afterwards so
  the cube stays cross-down; the solution remains valid, and the printed move list
  contains those rotations.
  \item \textbf{One puzzle.}  Only the $3\times3\times3$ is supported.  The cube
  model, the notation layer and the notation grammar are written for a cube of six
  faces and would need work for larger puzzles.
\end{enumerate}

\subsection{Verifying an installation}

\begin{lstlisting}[style=shell]
scripts/check.sh          # ruff format --check, ruff check, full pytest
uv run rubiks-cfop --scramble "R U R' U' F2 D" --no-color
\end{lstlisting}

The test suite covers the move tables (against Kociemba's published permutations),
the parser, the cubie invariants, both scramblers, the cross solver's optimality,
the F2L tables, exhaustive coverage of the last layer, the pipeline's stage order
and output format, and the command line itself.  A green run is the quickest proof
that the installation is sound.


\appendix
\section{A plan for the graphical interface}
\label{app:gui}

This appendix plans the desktop interface that sits on top of the engine.  The
engine was built headless precisely so that this work could be added without
rewriting it, and the plan below leans on what is already there.

\textbf{Status.}  Milestones M0 to M3 are implemented in the \texttt{rubiks\_cfop.gui}
package and start with \texttt{rubiks-cfop-gui} (or \texttt{python -m
rubiks\_cfop.gui}): the scramble bar with random scrambles and seeds, the net
canvas drawn through a painter, the solution table and a move-by-move player.  From
M4, a superseded job's result is ignored rather than the solve being canceled,
because \texttt{solve\_cube} does not take a cancellation token yet.  The plan for a
3D view is in \texttt{docs/gui-3d-plan.md}.

\subsection{Goals and non-goals}

\textbf{Goals.}

\begin{enumerate}[leftmargin=1.4em]
  \item Show the scrambled cube as a clear, colored net, and let the user read a
  solution stage by stage.
  \item Accept a scramble typed by the user, a facelet string from another tool, or
  a generated scramble with a reproducible seed.
  \item Let the user choose the cross color, or let the program choose the best one
  and say what it chose.
  \item Animate a solution move by move, with play, pause, step and a stage scrubber.
  \item Stay responsive: never block the window while solving, and always allow the
  solve to be canceled.
  \item Present errors as information, not as a crash: a bad token should be shown
  at its position in the input.
\end{enumerate}

\textbf{Non-goals.}

\begin{itemize}[leftmargin=1.4em]
  \item A photorealistic 3D cube.  A 2D net is unambiguous, cheap to draw and much
  easier to test; a 3D view is a later, optional renderer behind the same interface.
  \item Competition features: no timer, no inspection countdown, no session
  statistics, no scrambles from the network.
  \item Other puzzles, and any change to the solving algorithm.
\end{itemize}

\subsection{What the engine already provides}

Every one of these exists today and was designed for this purpose:

\begin{longtable}{@{}p{4.8cm}p{9.4cm}@{}}
\toprule
Engine facility & What the interface does with it \\
\midrule
\endhead
\texttt{Cube.apply\_scramble} & turns the typed scramble into a state, raising
\texttt{NotationError} with the token and index to highlight \\
\texttt{solve\_cube} & the whole CFOP solution, already verified, with the cross
color, the rotation applied and every stage \\
\texttt{Solution.to\_dict} & the payload the interface renders: stage labels, move
lists, and the facelets after each stage \\
\texttt{Stage.metadata} & the OLL number, the PLL name, the F2L slot and case, for
informative rows \\
\texttt{cross\_options} & the six crosses with their lengths, for the picker \\
\texttt{net\_layout} & a structured grid of colors; no escape codes to parse, and
no dependency on the terminal \\
\texttt{NetLayout.to\_dict} & the six faces as color names, ready for a web view \\
\texttt{DEFAULT\_PALETTE} & the standard colors, so the canvas matches the CLI \\
\texttt{CancelToken} & canceling a solve from the interface thread \\
\texttt{preload} & building the scrambler tables at start-up, off the UI thread \\
\texttt{RubiksError.to\_dict} & structured errors to display without parsing prose \\
\bottomrule
\end{longtable}

\subsection{Small engine additions the interface needs}

Four additions, all backwards compatible and all testable without a display:

\begin{enumerate}[leftmargin=1.4em]
  \item \texttt{iter\_move\_states(cube, moves)} --- an iterator over
  \texttt{(move, state)} pairs, so the animation applies one move at a time instead
  of owning a copy of the cube model.  Trivial to implement, and it keeps move
  application in one place.
  \item \texttt{Session} --- a stepwise solver: \texttt{next\_stage()} returns the
  next \texttt{Stage}, with \texttt{state}, \texttt{done} and \texttt{progress}
  properties.  A one-shot solve is then just a loop over it, so the interface can
  show a stage as soon as it is ready.
  \item \texttt{on\_stage} callback on the solve options, called from the solving
  thread as each stage completes.  Documented as \emph{called from that thread}:
  the interface posts to its own queue.
  \item \texttt{schemas/solution.schema.json} and a \texttt{schema\_version} field in
  the payload, plus a contract test that validates the emitted JSON against the
  schema.  This is what lets a second front-end (the web one) rely on the format.
\end{enumerate}

None of these adds a dependency to the engine, and none changes an existing
signature.

\subsection{Architecture}

The interface is split in three, and only the last layer knows about a toolkit:

\begin{lstlisting}[style=data]
+-------------------------------------------------------------+
|  app shell (main window, menu, shortcuts, settings dialog)  |
+-------------------------------------------------------------+
|  views        ScrambleBar  CubeCanvas  SolutionTable        |
|               Player       CrossPicker StatusBar            |
+-------------------------------------------------------------+
|  view-model   AppState: cube, solution, selected stage,     |
|               selected move, busy flag, last error          |
|               pure functions, no toolkit, fully unit tested |
+-------------------------------------------------------------+
|  engine       rubiks_cfop  (already headless, thread-safe)  |
+-------------------------------------------------------------+
\end{lstlisting}

Two rules make the plan work:

\begin{enumerate}[leftmargin=1.4em]
  \item \textbf{The view-model never imports a toolkit.}  It holds the state and the
  transitions --- \emph{start a solve}, \emph{select a stage}, \emph{step a move},
  \emph{report an error} --- and is tested with plain pytest.  Everything that is
  easy to get wrong lives here, where it can be tested headlessly.
  \item \textbf{The canvas draws through a painter interface.}  \texttt{CubeCanvas}
  receives a \texttt{Painter} with \texttt{fill\_rect(x, y, w, h, color)} and
  \texttt{draw\_text(...)}.  Production uses the toolkit's painter, tests use a
  recorder, so the layout can be asserted without a screen.
\end{enumerate}

\subsection{Choosing a toolkit}

\begin{longtable}{@{}p{3cm}p{2.4cm}p{4.2cm}p{4.4cm}@{}}
\toprule
Option & Dependency & Strengths & Costs \\
\midrule
\endhead
Tkinter & none (standard library) & matches the engine's zero-dependency stance;
ships with Python; enough for a net, a table and a canvas & dated look; the
\texttt{Canvas} widget is a little low level \\
PySide6 / Qt & a large wheel & polished widgets, good high-DPI support, animation
timers & a heavy dependency and a license to respect; overkill for a net and a
table \\
Web front-end & a server plus a browser & the JSON payload already suits it;
portable; easy to share & a second process, a second language, and no offline
single-file story \\
\bottomrule
\end{longtable}

\textbf{Recommendation:} Tkinter for the first interface.  It keeps the project
installable with no optional dependency, and the view-model plus painter split means
a Qt or web front-end can replace the view layer later without touching the
reasoning.  A web front-end is the natural second step, and the JSON payload and its
schema are the contract for it.

\subsection{Screens and widgets}

\begin{description}[leftmargin=0em,style=nextline]
  \item[Scramble bar] one line of text in WCA notation, with a \emph{Random}
  button, a seed field and a \emph{lock seed} toggle, a style selector
  (random-state, or random-move with a warning that it is not WCA-legal), and a
  \emph{Solve} button.  Invalid notation is marked at the offending token, taken
  from \texttt{NotationError.index}.
  \item[Cube canvas] the unfolded net, drawn from \texttt{net\_layout} in the chosen
  palette.  It shows either the scramble or the state after the selected stage, and
  it is the surface the animation draws on.
  \item[Solution table] one row per stage: label, case (the OLL number, the PLL name,
  the F2L slot), the moves, the length and the running total.  Selecting a row moves
  the canvas to that stage's \texttt{facelets} and the player to that stage.
  \item[Player] play, pause, step forward and back, a speed slider and a stage
  scrubber.  Playing a stage animates its moves one at a time; stepping shows the
  cube after a single move.
  \item[Cross picker] the six colors with their move counts, best first, from
  \texttt{cross\_options}; picking one re-solves with that cross and reports where
  the cube was rotated.
  \item[Status bar] the version, the time spent, the total move count and a
  \emph{verified} badge that comes from the payload, plus a cancel button while a
  solve runs.
  \item[Settings] palette (standard or a high-contrast set), animation speed, font
  size, and whether a generated scramble is remembered between sessions.
\end{description}

\subsection{Threading, responsiveness and cancellation}

\begin{enumerate}[leftmargin=1.4em]
  \item \textbf{Nothing expensive on the UI thread.}  A CFOP solve takes up to a few
  seconds and building the scrambler tables takes about ten.  Both run on worker
  threads.
  \item \textbf{One solve at a time.}  Starting a second solve cancels the first;
  the token makes that prompt and the \texttt{CancelledError} is reported as
  \emph{canceled}, not as a failure.
  \item \textbf{Results cross threads as data.}  The worker hands back
  \texttt{Solution.to\_dict()}; the UI thread rebuilds the view-model from it.  No
  engine object is shared while it is being computed.
  \item \textbf{Table preloading at start-up.}  A splash or a status message during
  \texttt{preload()}, which then makes every later scramble immediate.
  \item \textbf{Cooperative timers for animation.}  The player advances one move per
  tick and stops when the window is not visible; the engine is only asked for the
  next state, never for a whole animation.
\end{enumerate}

\subsection{Testing without a display}

A graphical interface is usually the least tested part of a project.  The
architecture above makes most of it testable anyway:

\begin{longtable}{@{}p{4.6cm}p{9.4cm}@{}}
\toprule
Layer & How it is tested \\
\midrule
\endhead
view-model & plain pytest: transitions, stage selection, move stepping, error
reporting, busy and canceled states \\
canvas layout & a recording painter: the sequence of rectangles and colors must
match \texttt{net\_layout} cell for cell \\
payload contract & the schema file plus a test that validates
\texttt{Solution.to\_dict()} and the CLI's \texttt{--json} against it \\
engine & the existing suite; the interface adds no logic below the view-model \\
interactive smoke test & a scripted session on a headless display where one is
available, skipped otherwise, that types a scramble, solves, steps twice and
asserts the canvas state \\
\bottomrule
\end{longtable}

The rule of thumb: if a behavior cannot be tested without a screen, it belongs in
the view-model or below, not in the widget.

\subsection{Milestones}

Each milestone leaves the program in a usable state and ends with a green gate
(formattingl, linting and the whole test suite).

\begin{longtable}{@{}p{1.5cm}p{5.4cm}p{7.2cm}@{}}
\toprule
Step & Scope & Acceptance \\
\midrule
\endhead
M0 & Skeleton: window, view-model, canvas with the painter interface, static net &
\texttt{rubiks-cfop-gui --scramble "R U R' U'"} shows the net; view-model and canvas
layout tests green; no change to the engine \\
M1 & Scramble bar, solve button, solution table, verified badge & twenty random
scrambles solved through the interface; the payload equals the CLI's
\texttt{--json}; a bad token is reported at its index \\
M2 & Player and animation & play, pause, step and speed animate a solution;
stepping lands on the exact \texttt{facelets} of each stage \\
M3 & Cross picker, seed field, preload splash & choosing a color re-solves with
it; the same seed reproduces the same scramble; no freeze during preload \\
M4 & Cancellation, progress, structured error dialogs & cancel returns within a
fraction of a second and leaves the window usable; every error class renders \\
M5 & Settings, palette, packaging, documentation & a \texttt{[gui]} extra and a
\texttt{rubiks-cfop-gui} entry point; screenshots in the manual; the suite is still
green without a display \\
\bottomrule
\end{longtable}

\subsection{Packaging}

\begin{lstlisting}[style=shell]
# optional dependency group, so the engine stays dependency free
uv add --optional gui pytest-qt   # or nothing at all, for Tkinter

# entry point
rubiks-cfop-gui --scramble "R U R' U'" --cross white
\end{lstlisting}

With Tkinter the group is empty, which is the strongest argument for starting there:
\texttt{pip install rubiks-cfop} gives a working command line, and
\texttt{pip install 'rubiks-cfop[gui]'} gives the window.

\subsection{Risks}

\begin{longtable}{@{}p{5.4cm}p{8.8cm}@{}}
\toprule
Risk & Mitigation \\
\midrule
\endhead
The toolkit choice turns out to be wrong & the view-model and the painter interface
isolate it; a replacement touches only the view layer \\
Slow solves make the interface feel broken & the two F2L improvements in
Section~\ref{sec:trouble} are the real fix; meanwhile the progress callback and the
cancel button keep it honest \\
Animation looks choppy & redrawing 54 cells per frame is negligible; if needed,
animate only the affected layer \\
The engine API drifts while the interface is written & the payload schema and its
contract test are the interface between them \\
Nothing can be tested on the build machine & the canvas painter and the view-model
need no display; the interactive test is skipped when there is no display \\
\bottomrule
\end{longtable}

\subsection{Order of work}

M0 and M1 are deliberately small: they prove the two risky assumptions --- that the
canvas can be tested headlessly, and that a solve can be run off the UI thread
without freezing --- before any polish.  If those hold, M2 to M5 are ordinary
interface work, and the web front-end becomes a matter of writing a second view over
the same payload.


\end{document}
